% --- Άσκηση 35385 (35385.tex) ---
\Askhsh{35385}{4}
\begin{ylh}{2}
    \item 2.1. Οι Πράξεις και οι Ιδιότητές τους
    \item 2.4. Ρίζες Πραγματικών Αριθμών
    \item 3.3. Εξισώσεις 2ου Βαθμού
    \item 4.2. Ανισώσεις 2ου Βαθμού
    \item 6.1. Η Έννοια της Συνάρτησης
    \item 6.2. Γραφική Παράσταση Συνάρτησης
\end{ylh}
Δίνονται οι συναρτήσεις $f(x) = x^{2} + \beta$,$\ \ g(x) = x + \beta$, όπου $x \in R$ και $\beta$ σταθερός πραγματικός αριθμός. Είναι γνωστό ότι η γραφική παράσταση της $g(x)$ διέρχεται από το σημείο $M\left( \dfrac{3\beta}{2},\  - 3 - \dfrac{\beta}{2} \right).$
\begin{alist}
\item Να αποδείξτε ότι $\beta = \mathbf{-}\ 1$.
\monades{6}
\item Για $\beta = \mathbf{-}\ 1$
\begin{rlist}
\item
  Να βρείτε τα σημεία τομής της γραφικής παράστασης της συνάρτησης $f(x)$ με τους άξονες $x΄x$, $y΄y.$
\monades{5}
\item
  Να βρείτε τα διαστήματα στα οποία η γραφική παράσταση της $f(x)$ βρίσκεται κάτω από τη γραφική παράσταση της συνάρτησης $g(x).$
\monades{7}
\item
  Να λύσετε την εξίσωση $\dfrac{f(x)}{g(x)} + \dfrac{g(x)}{f(x)} = 3$.
\monades{7}
\end{rlist}
\end{alist}
